\section{Reconstruction}
\label{sec:reconstruction}

The fill rules of Section~\ref{sec:fillmodels} need four things the dataset does
not state directly: the block each event belongs to, the volume resting ahead of
each order, the moment each order reaches the front of its queue, and which side
of each trade rested. We summarise how each is obtained; details and validation
are in Appendix~\ref{app:reconstruction}.

\paragraph{Blocks.}
The dataset carries no block number, but status timestamps are block times:
96--98\% of them coincide exactly between \BTC{} and \SOL, against $4\times10^{-4}$
chance coincidences for event times. A block is recovered by grouping events on
equal timestamps. Within a block, \textbf{cancellations precede fills in
98.8--99.4\% of event pairs}, so a cancellation submitted in the same block as an
aggressor is processed first (Section~\ref{sec:block}).

\paragraph{Queue position and the front of the queue.}
\label{sec:queue}
The book-diff stream reports the size each order retains after every change, so
the volume of a level is a running sum of size changes and the volume ahead of
an order at placement is the level's volume just before it arrives. The moment
an order reaches the front---when nothing that arrived earlier at its level is
still resting---is found by a linear sequential pass per level. Size reductions
keep priority. Book diffs carry no time; each change takes the time of the
corresponding order status.

\paragraph{Fills.}
A fill is taken from the trade stream, the maker being the side whose order is in
the book. Fill statuses cannot be used: they record only about 60\% of maker
fills, and the omission is one-sided, so it would bias outcomes rather than add
noise (Appendix~\ref{app:reconstruction}).

\paragraph{Validation.}
Every rebuilt hour passes checks that count states impossible in a correctly
ordered book---negative level volume, a crossed book, a carried size disagreeing
with the size recorded in the event---all zero on real data
(Table~\ref{tab:checks}).
